两个检查点的129条序列均已完成，每点执行1,032次网络前向；
184份官方充分状态通过哈希核验，40份总体及当前帧指标独立复算一致。
表\ref{tab:trained-memory}显示，在完全共享局部候选的条件下，持久记忆改善了任务AP和gap恢复召回。
E0下的AP提高2.5343个百分点，E1下提高2.2785个百分点；两者仍低于相应完整前缀方式。
完整的$T=2$至$T=5$结果、当前帧指标和事件分母保留于随附CSV。

\begin{table}[t]
\centering
\small
\setlength{\tabcolsep}{3pt}
\caption{固定感知器的五复访记忆对照。AP和恢复召回为百分数。
每行129条序列、6个物理参考场景；恢复计数为正确恢复次数/1,101个gap机会，
包含同场景不同序列的重复事件，不是1,101个独立样本。}
\label{tab:trained-memory}
\begin{tabular}{llrrr}
\toprule
Model & Method & AP & Recall & Correct/gaps\\
\midrule
E0 & FH-B2 & 4.5902 & 4.2688 & 47/1101\\
E0 & LOCAL-B2 & 1.3417 & 4.7230 & 52/1101\\
E0 & LOCAL-B4 & 3.8760 & 20.5268 & 226/1101\\
\midrule
E1 & FH-B2 & 5.9381 & 6.6303 & 73/1101\\
E1 & LOCAL-B2 & 1.1343 & 6.5395 & 72/1101\\
E1 & LOCAL-B4 & 3.4128 & 26.4305 & 291/1101\\
\bottomrule
\end{tabular}
\end{table}

LOCAL-B4相对LOCAL-B2的恢复召回差在E0下为15.8038个百分点，
六场景配对bootstrap的95\%区间为[12.6850,23.0576]；
E1下为19.8910个百分点，区间[16.1261,28.4120]。
采用10,000次参考场景重采样，区间描述该固定面板的事件率不确定性，不能替代独立训练重复。
身份切换率差分别为$-1.3858$和$-0.0470$个百分点，
区间分别为[$-3.6842$,0.1945]和[$-2.2364$,2.4506]，均跨零。
因此不能将恢复召回改善推广为已确认的身份切换率下降。
E1下的局部持久状态AP反而低于E0下的3.8760\%；
完整前缀感知指标的提升未自动转化为固定记忆配置下的任务AP提升。

五复访时，E0的完整前缀与局部持久方式的网络加关联耗时中位数分别为969.30和529.40毫秒，
E1对应为965.14和526.96毫秒。两种权重的完整/局部网络峰值分配显存中位数
分别为3.0910/1.6163 GiB。时间不包含单列的输入搬运和后处理，不能称为端到端总延迟；
两个局部方式复用同一感知前向。
显式持久状态为61,008字节，E0/E1最大占用分别为38/40个槽，总容量100，均无拒绝新实体事件。
所测序列最长五次复访且没有容量饱和，尚不能验证更长历史下的淘汰策略或信息保留能力。
